\begin{enumerate}
\item Tìm tọa độ tâm $I(a; b)$ và bán kính $R$.
\item Thay vào dạng $(x-a)^2 + (y-b)^2 = R^2$[cite: 1].
\item Nếu bài cho đi qua $3$ điểm, gọi dạng $x^2 + y^2 - 2ax - 2by + c = 0$[cite: 1], thay tọa độ $3$ điểm vào lập hệ $3$ phương trình ẩn $a, b, c$ và bấm máy giải hệ.
\end{enumerate}
Lưu ý: Đường tròn đường kính $AB$ có tâm là trung điểm $AB$ và $R = \frac{AB}{2}$[cite: 1]. Đường tròn tâm $I$ tiếp xúc đường thẳng $\Delta$ thì $R = d(I, \Delta)$[cite: 1]. Đừng quên bình phương $R$ khi viết phương trình[cite: 1].
